% sec:fd-truncation — Truncation Error Analysis
\section{Truncation Error Analysis}
\label{sec:fd-truncation}

The stencils of \cref{sec:fd-stencils,sec:fd-orders} carry a promise:
refine the grid and the approximation improves at a predictable rate.  But
the promise is only useful if we can quantify the error --- both to choose
a resolution that meets accuracy requirements and to verify, after the
fact, that the code produces the expected convergence.  The truncation
error is the formal object that makes both tasks precise.

\paragraph{Local truncation error.}

Applying a stencil to exact function values produces a derivative
approximation that differs from the true derivative by a remainder.  This
remainder is the \emph{local truncation error}: the error committed at a
single grid point on a single evaluation, assuming the underlying function
is known exactly.

To define it, consider a stencil with weights $\{w_{k_j}\}$ and offsets
$\{k_j\}$ that approximates $f^{(p)}_i$.  The stencil formula is
%
\begin{equation}
  D_p[f]_i \;\equiv\;
  \frac{1}{\Delta x^p}\,\sum_{j=1}^{N} w_{k_j}\,f_{i+k_j} \,,
  \label{eq:stencil-operator}
\end{equation}
%
where the weights satisfy the Vandermonde conditions
(\cref{eq:stencil-conditions}).  The local truncation error $\tau_i$ is
the difference between this discrete operator and the exact derivative:
%
\begin{equation}
  \tau_i \;\equiv\; D_p[f]_i - f^{(p)}_i \,.
  \label{eq:lte-def}
\end{equation}
%
This is a pointwise quantity: it depends on the grid spacing, the stencil
order, and the local smoothness of $f$.

\paragraph{Structure of the error.}

Substituting the Taylor expansion (\cref{eq:taylor-grid}) into
\cref{eq:stencil-operator} and using the Vandermonde conditions to cancel
the first $N$ terms gives a series in even powers of $\Delta x$
for centred stencils.  The leading term has the form
%
\begin{equation}
  \tau_i = C_q\,\Delta x^q\,f^{(p+q)}_i
           + \order{\Delta x^{q+2}} \,,
  \label{eq:lte-leading}
\end{equation}
%
where $q = 2m$ is the order of accuracy (the stencil uses $2m + 1$ points)
and $C_q$ is a rational coefficient determined by the uncancelled term in
the Vandermonde expansion.  The integer $p + q$ is the order of the
derivative that controls the error magnitude: higher-order stencils push
this to higher derivatives of $f$, which are typically smaller for smooth
fields.
\physnote{The error depends on derivatives of $f$ that were not part of
the original approximation.  If $f$ has large high-order derivatives ---
as near discontinuities or sharp features --- the truncation error can be
much larger than the leading-term estimate suggests.}

The error coefficients tabulated in \cref{sec:fd-orders} are specific
values of $\abs{C_q}$: for the fourth-order first derivative,
$\abs{C_4} = 1/30$ and the leading error magnitude is
$(\Delta x^4/30)\,\abs{f^{(5)}_i}$; for the fourth-order second
derivative, $\abs{C_4} = 1/90$ and the leading error magnitude is
$(\Delta x^4/90)\,\abs{f^{(6)}_i}$.  These coefficients are not
arbitrary --- they are fixed by the Vandermonde system and can be computed
for any stencil.

\paragraph{Order of accuracy and convergence rate.}

The exponent $q$ in \cref{eq:lte-leading} controls the convergence rate
under grid refinement.  Halving the grid spacing replaces $\Delta x$ with
$\Delta x/2$, so the truncation error decreases by a factor of $2^q$:
%
\begin{equation}
  \frac{\tau_i(\Delta x)}{\tau_i(\Delta x/2)}
  = 2^q + \order{\Delta x^2} \,,
  \label{eq:convergence-ratio}
\end{equation}
%
provided $f^{(p+q)}_i \neq 0$ (if the leading derivative vanishes, the
ratio is dominated by the next term and may differ from $2^q$).
For a second-order stencil ($q = 2$), halving the grid reduces the error
by a factor of $4$; for a fourth-order stencil ($q = 4$), the factor is
$16$.  This ratio is directly measurable: compute a stencil derivative at
two resolutions, compare both to the exact answer (or to a much
finer-resolution computation), and check that the error ratio matches $2^q$.
This is the core idea behind the convergence tests of
\cref{sec:fd-convergence}.
\implnote{The test \texttt{test\_d1x\_sin\_convergence} in
\codemodule{stencils} checks that the error ratio between successive
doublings of $N$ lies in $[12, 20]$, bracketing the fourth-order
prediction of $16$.}

\paragraph{From local to global error.}

The truncation error $\tau_i$ characterises the error at a single point.
When a finite difference scheme is used to solve a differential equation
--- stepping forward in time using spatial derivatives --- the errors from
every grid point and every time step accumulate.  The relationship between
the local truncation error and the global solution error depends on the
stability of the scheme, a topic treated in the ODE integration chapter
(\cref{ch:ode-integration}).  The essential result is the
\emph{Lax equivalence theorem}: for a consistent, stable linear scheme,
the global error converges at the same order as the local truncation error
\cite{LeVeque2007}.
\xrefnote{Stability analysis and the CFL condition are developed in
\cref{ch:ode-integration}.  The Lax equivalence theorem links the
truncation error analysis here to the global convergence of the full
discretisation.}

Consistency requires that the truncation error vanish as $\Delta x \to 0$,
which is guaranteed by construction for any stencil derived from the
Vandermonde system.  Stability requires that the time-stepping scheme does
not amplify errors unboundedly --- this is where the CFL condition and
Kreiss--Oliger dissipation (\cref{sec:fd-dissipation}) enter.  When both
conditions hold, the global error inherits the $\order{\Delta x^q}$
scaling of the local truncation error, and the convergence ratio
\cref{eq:convergence-ratio} applies to the solution as a whole.  In
practice, this means that the local convergence tests we can run on
isolated stencil operators predict the accuracy of the full evolution
--- the most consequential result in this chapter.
\warnnote{The Lax equivalence theorem applies to linear problems.
Nonlinear systems --- including the BSSN equations --- require additional
analysis, but the convergence rate $\order{\Delta x^q}$ is observed
empirically and confirmed by the self-convergence tests of
\cref{sec:fd-convergence}.}

\paragraph{Richardson extrapolation.}

The structure of the truncation error (\cref{eq:lte-leading}) implies that
the dominant error term is a known function of $\Delta x$.  This opens a
path to higher accuracy without higher-order stencils: compute the same
derivative at two different resolutions and combine the results to cancel
the leading error term.

Suppose a quantity $Q$ is computed at grid spacings $\Delta x$ and
$\Delta x/2$, with errors dominated by the leading truncation term:
%
\begin{align}
  Q(\Delta x)   &= Q_{\text{exact}} + C\,\Delta x^q
                   + \order{\Delta x^{q+2}} \,, \label{eq:richardson-coarse} \\
  Q(\Delta x/2) &= Q_{\text{exact}} + C\,\frac{\Delta x^q}{2^q}
                   + \order{\Delta x^{q+2}} \,.
                   \label{eq:richardson-fine}
\end{align}
%
The coefficient $C$ is the same in both expressions --- it depends on the
derivatives of $f$, not on $\Delta x$ --- so we can eliminate it.
Multiplying \cref{eq:richardson-fine} by $2^q$ and subtracting
\cref{eq:richardson-coarse} gives the Richardson extrapolant:
%
\begin{equation}
  Q_R \;=\; \frac{2^q\,Q(\Delta x/2) - Q(\Delta x)}{2^q - 1}
  \;=\; Q_{\text{exact}} + \order{\Delta x^{q+2}} \,.
  \label{eq:richardson-extrapolant}
\end{equation}
%
The extrapolant $Q_R$ has an error of order $\Delta x^{q+2}$ --- two
orders higher than either input.  A fourth-order stencil ($q = 4$)
combined with Richardson extrapolation yields sixth-order accuracy, at the
cost of evaluating the stencil at two resolutions.
\histnote{Richardson introduced this technique in 1911 for improving the
accuracy of numerical integration \cite{Richardson1911}.  It is
sometimes called ``deferred approach to the limit'' and underpins Romberg
integration and the Bulirsch--Stoer ODE method.}

\paragraph{Richardson extrapolation as a diagnostic.}

In practice, the extrapolation is more valuable as a verification
tool than as an accuracy booster.  The key observation is that the
extrapolant (\cref{eq:richardson-extrapolant}) is valid only when the
error is dominated by the leading truncation term --- that is, when the
grid is fine enough that higher-order terms are negligible.  If the
extrapolant and the fine-grid value agree to high precision, the
computation is in the asymptotic regime and the convergence rate is
trustworthy.  If they disagree significantly, the grid is too coarse for
the leading-term approximation to hold, and the reported convergence rate
may be misleading.

This diagnostic role is especially important for the BSSN equations, where
gauge dynamics, constraint violations, and boundary effects can all
contaminate the error.  A Richardson extrapolation that yields the
expected order increase confirms that truncation error --- not some
other source --- dominates the solution error at the resolutions being
tested.  The convergence-testing framework of \cref{sec:fd-convergence}
uses this principle systematically.

\paragraph{A worked example.}

To make the extrapolation concrete, return to the test function
$f(x) = \sin(2\pi x)$ from \cref{sec:fd-orders}.  The fourth-order first
derivative at $x = 0.5$ has errors of approximately
$2.04 \times 10^{-3}$ at $N_x = 20$ and $1.3 \times 10^{-4}$ at
$N_x = 40$ (\cref{eq:d1-error-fourth-order}).  With $q = 4$, the
Richardson extrapolant is
%
\begin{equation}
  Q_R = \frac{16\,Q(\Delta x/2) - Q(\Delta x)}{15} \,.
  \label{eq:richardson-worked}
\end{equation}
%
Denoting the errors at each resolution by $\epsilon_1 \approx 2.04 \times
10^{-3}$ and $\epsilon_2 \approx 1.3 \times 10^{-4}$, the Richardson
residual is
$(16\,\epsilon_2 - \epsilon_1)/15
 = (2.08 \times 10^{-3} - 2.04 \times 10^{-3})/15
 \approx 3 \times 10^{-6}$
--- nearly two orders of magnitude smaller than the fine-grid error alone.
The improvement confirms that the error at these resolutions is indeed
dominated by the fourth-order truncation term; the residual is the
$\order{\Delta x^6}$ remainder that Richardson extrapolation cannot
cancel.
